\chapter{Structure Before Decoration}

A production system becomes useful when repeated editorial decisions are represented as stable rules rather than remembered adjustments. This synthetic chapter contains ordinary prose only. It is deliberately generic and exists to test the baseline book model.

The first paragraph after a chapter opening is allowed to follow the class convention. Later paragraphs demonstrate ordinary paragraph indentation and line breaking. A system should make such behavior predictable across titles without requiring page-by-page manual intervention.

\section{A Small Structural Test}

Reliable production separates canonical content from generated output. A PDF is an artifact of the production process, not a substitute for the source from which it was built. When the generated files are removed, the source should be sufficient to reproduce the book.

\subsection{Notes and emphasis}

This sentence contains \emph{emphasis}, a short quotation---``a reproducible book is easier to maintain''---and a footnote.\footnote{This synthetic note exercises the baseline note infrastructure without prescribing a final scholarly-note design.}

A visual baseline should also contain enough ordinary prose to reveal page rhythm rather than leaving a nearly empty continuation page. This paragraph therefore extends the synthetic fixture without introducing project-specific content or design assumptions.

The production rule remains structural: content should flow according to the profile, while exceptional manual page breaks are reserved for editorially meaningful cases rather than used as routine repairs for local spacing.
